\scanpage{88}

qui s'explicitent en termes de (78) en notant que tout drapeau de $J\cup L$ \inserted{intervenant dans (78)} s'écrit de façon unique
\[
(80)\qquad(j_{\unsure{1|q}}\mathord{\cdot}\mathord{-}\mathord{\cdot}j_p\ l_1\mathord{\cdot}\mathord{-}l_q)\qquad
\begin{gathered}[t]
i_{\unsure{q|0}}\mathord{-}\mathord{-}i_p\in J\\
l_1\mathord{-}\mathord{-}l_q\in L\\
p,q\geq1
\end{gathered}
\]
et \deleted{\unsure{ensuite}} on a des
\diagram{\begin{tikzcd}[ampersand replacement=\&,column sep=large,row sep=large]
(81)\&\Sigma_{j_1\mathord{-}\mathord{-}j_p\,l_1\mathord{-}\mathord{-}l_q}\arrow[r,hook]\arrow[d]\&\Sigma_{l_1\mathord{\cdot}\mathord{\cdot}\mathord{-}l_q}\\
\&\Sigma_{j_1,\mathord{\cdot}\mathord{-},j_p}\&
\end{tikzcd}}

\uline{Remarque} Les $\Sigma_{J,L}$ ont un sens \inserted{(d.s une catégorie isotopique convenable)} indépendamment des rétractions coniques locales autour des $X_i^*$, du moins en théorie topologique modérée, ainsi que les $\Sigma_{J,L}\hookrightarrow\Sigma_L$ --- et la description que j'en ai donné ici est adaptée à ce cas --- la structure $\Sigma_{J,L}\to\Sigma_J$ par contre utilise les \deleted{structures} \inserted{fibrations} \inserted{\deleted{\unsure{de ri.}}} \deleted{côniques} locales. Quand on dispose $=$ des structures côniques locales (ce qui, en topologie modérée, et dans le cas équisingulier, est une donnée supplémentaire ``anodine''), il est plus commode,
\margin{87}
